\subsubsection*{定理 1.3.}
设 $X$ 是特征指数为 $p$ 的域 $k$ 上维数 $n\geq1$ 的不可约完全交.
对每个整数 $d\geq n+1$, $X$ 的充分一般的 $d$ 次超曲面截面在 1.0 的意义下非特殊.

\medskip
\noindent\emph{证明.}
$X$ 的 $d$ 次超曲面截面 $X.H_u$ 由 $H^0(X,\mathcal O_X(d))$ 中的直线参数化,
亦即由下列空间参数化
\begin{equation}\tag{1.3.1}
U=\mathbb P\bigl(H^0(X,\mathcal O_X(d))\bigr).
\end{equation}
在 $U$ 上有“$X$ 的 $d$ 次超曲面泛截面”,
\begin{equation}\tag{1.3.2}
f:X.H_u\longrightarrow U.
\end{equation}
设 $U'$ 是由使 $\dim(X_u)=n-1$ 的 $u\in U$ 组成的开子集.
由 XI~1.5(i) 及 2.5, 限制
\begin{equation}\tag{1.3.3}
\left.R^{n-1}f_*\mathcal O_{X.H_u}\right|_{U'}
\end{equation}
局部自由, 其形成与任意基变换可交换, 且其秩 $M$ 至少为一.
若 $p=1$, 这就对每个次数至少为 $n+1$ 的超曲面截面证明了定理.

因此设 $p>0$. 迭代 Hasse--Witt 运算 $\mathfrak F_{p^M}$ 同样局部由元素属于
$\mathcal O_U$ 的 $M\times M$ 矩阵表示. 因此,
\begin{equation}\tag{1.3.4}
\{u\in U'\mid X.H_u\text{ 非特殊}\}
\quad\text{是下述空间的开集: }U'.
\end{equation}
只需证明 (1.3.4) 中的集合非空.

先考虑情形
\begin{equation}\tag{1.3.5}
k=\mathbb F_q.
\end{equation}
将 $\mathbb F_q$ 换成有限扩张后, 可设
\begin{equation}\tag{1.3.6}
\operatorname{card}\bigl(X(\mathbb F_q)\bigr)>0
\end{equation}
且
\begin{equation}\tag{1.3.7}
\begin{aligned}
&\text{存在齐次坐标}\\
&\qquad X_1,\dots,X_{n+r+1}\text{ 位于环境空间 }\mathbb P_k^{n+r}\\
&\text{使 }X\text{ 不交于如下线性子簇}\\
&\qquad X_1=\cdots=X_{n+1}=0.
\end{aligned}
\end{equation}
选取次数为 $d$ 的齐次式
\begin{equation}\tag{1.3.8}
H(X_1,\dots,X_{n+1})\in
\mathbb F_q[X_1,\dots,X_{n+1}],
\end{equation}
使其具有性质 (1.2.2); 由于 $d\geq n+1$, 这是可能的.
也用 $H$ 表示 $\mathbb P^{n+r}$ 中由 $H(X_1,\dots,X_{n+1})=0$ 定义的超曲面.
由构造, $X.H$ 没有 $\mathbb F_q$ 有理点. 根据 (1.3.6), 这也表明 $X\not\subset H$;
由于 $X$ 不可约, $X.H$ 维数为 $n-1$, 故 $H\in U'$.
现在命题~1.1.5 表明 $X.H$ 非特殊.

\medskip
\noindent\emph{底本注.}
照录法文层保留了印本的指数 $p^n$、底层 $\mathcal O_u$ 以及简写引用“(1.5)”.
秩 $M$ 的幂零性判别、基底 $U$ 以及命题~1.1.5 要求采用上文相容的读法.
(1.3.1) 中的右括号仅在继承的 TeX 中缺失, 已直接按扫描件补回.
